\title{LongRoPE: Extending LLM Context Window Beyond 2 Million Tokens}

\begin{abstract}
The ability of large language models (LLMs) to handle extensive context windows is highly sought after. Nevertheless, present methods for enlarging context windows face limitations—mainly due to the significant costs associated with fine-tuning, limited availability of lengthy textual data, and problematic extreme values that arise from newly introduced token positions—resulting in current extended context windows being capped at approximately 128k tokens.
\end{abstract}

This work presents LongRoPE, a method that, for the first time, expands the context window of pre-trained LLMs to an impressive \textbf{2048k} tokens, utilizing as few as 1k fine-tuning steps with training sequences up to 256k tokens, all while preserving the models' capability on their original, shorter context window. This advancement is made possible by three primary contributions: (i) the identification and utilization of two types of non-uniformities in positional interpolation, discovered through an efficient search mechanism, which leads to improved fine-tuning initialization and enables an eightfold expansion without further fine-tuning; (ii) the development of a staged extension approach, where a 256k-token LLM is first fine-tuned, followed by a secondary positional interpolation on the extended model to realize the 2048k-token window; (iii) refinement of LongRoPE at an 8k-token window to ensure performance on shorter contexts is fully recovered. Comprehensive evaluations on LLaMA2 and Mistral spanning a range of tasks confirm the robustness of our approach. Models modified with LongRoPE retain their original structure aside from slight alterations to the positional embedding component and can still leverage most previous optimizations. The implementation will be made available at \texttt{https://github.com/microsoft/LongRoPE}

\section{Introduction}

Although Large Language Models (LLMs) have achieved substantial advancements across a range of tasks~\cite{gpt4,llama2}, their context window size remains constrained, such as the 4096-token boundary in LLaMA2~\cite{llama2}. When input sequences exceed this limit, LLMs exhibit decreased performance because they encounter token positions not present during their training phase. This limitation creates difficulties in applications that demand extended context, including in-context learning with many instances~\cite{Huang2023FewerIM} and LLM agents~\cite{park2023generative,madaan2023selfrefine}.

Recent studies demonstrate that the context window of a pre-trained LLM can be enlarged to approximately 128k tokens through fine-tuning with longer sequences~\cite{longlora,pi,yarn,zhang2024soaring,liu2023scaling}. However, pushing the limits of the context window faces three primary challenges. Firstly, introducing novel position indices that have not been previously trained leads to numerous anomalous values, creating out-of-distribution problems and hindering the convergence of the fine-tuning process~\cite{pi}. This difficulty is exacerbated when expanding the context from 4k to over 1000k tokens, as this results in more than 90\% of the positions being newly generated. Secondly, successful fine-tuning requires training materials that match the extended lengths, yet there is a scarcity of datasets containing texts longer than 1000k tokens. Additionally, processing and training on such massive sequences demands excessive computational resources and lengthy training times, posing substantial logistical barriers. Finally, with extremely extensive context windows, the attention mechanism suffers from dilution, as the model distributes attention across an overwhelming number of tokens, thereby diminishing performance on shorter contexts~\cite{pi}.

A common strategy to address the initial issue involves the interpolation of RoPE positional embeddings~\cite{rope,pi}, wherein newly introduced position indices are rescaled to match the range encountered during pretraining, illustrated in Fig.\ref{fig:overview}. The Position Interpolation (PI) method~\cite{pi} achieves this by performing linear interpolation on RoPE's rotary angles according to the expansion ratio. NTK~\cite{ntk,dynamicntk} recommends non-uniform interpolation and extrapolation along different RoPE dimensions. Meanwhile, YaRN~\cite{yarn} partitions the RoPE dimensions into three groups based on their frequency characteristics, subsequently implementing extrapolation, NTK interpolation, and linear interpolation for each group. Nevertheless, the positional embeddings in the Transformer model inherently possess \emph{complex, non-uniform information entropy}. The current methods fail to adequately exploit these intricate non-uniform patterns, resulting in lost information and thus constraining the size of the achievable context window.

Section~\ref{sec:analysis} presents two main empirical insights: \textbf{(1)} For positional interpolation to be effective, it is crucial to account for two types of non-uniformities—those arising from differences across RoPE dimensions and those tied to specific token positions. Minimal interpolation is preferable for lower RoPE dimensions and tokens at the beginning of sequences; however, the precise configuration is contingent upon the intended extrapolation length. \textbf{(2)} Incorporating these sources of non-uniformity into the interpolation process enhances the preservation of information encoded by the original RoPE, especially within essential dimensions and at critical token locations. As a result, positional information loss is reduced, thereby facilitating more robust parameter initialization for subsequent fine-tuning. This approach also supports an extension of up to 8$\times$ without the need for fine-tuning.

Building on these results, we propose \textbf{LongRoPE}, a robust approach that extends the LLM context window to over 2 \emph{million} tokens. The design of LongRoPE relies on three principal components. Firstly, {LongRoPE} leverages multidimensional variations in positional interpolation to their fullest potential. Specifically, it determines optimal rescale factors for the rotation angles in each RoPE dimension as a function of token positions. Given that the search space for suitable rescale factors grows exponentially with the increase in the context window, {LongRoPE} employs an evolutionary search strategy. This strategy incorporates two optimization methods to enhance the efficiency of the search process. An illustration of the optimized, rescaled RoPE generated by this method is presented in Fig.~\ref{fig:overview}.

Subsequently, {LongRoPE} implements a resource-efficient, incremental extension method to obtain a 2048k context window, circumventing the necessity for direct fine-tuning on exceedingly long input sequences, which are seldom found in practice. The approach initiates by identifying a 256k context length on the original pre-trained LLM, followed by fine-tuning the model at this specific length. Next, utilizing our non-uniform positional interpolation—enabling an 8$\times$ extension even when fine-tuning is not performed—we proceed with a secondary search for optimal RoPE rescale factors on the already fine-tuned, length-extended LLM. Through this process, we successfully reach the 2048k context window on models such as LLaMA2 and Mistral~\cite{mistral}.

To address potential declines in performance with the default (smaller) context length, {LongRoPE} maintains adjustments to the RoPE rescale factors even in the expanded LLM configuration. Employing a comparable strategy as in the transition from 256k to 2048k windows, we additionally apply downward scaling for 4k and 8k context windows on the LLM fine-tuned at 256k, using our search algorithm to minimize positional interpolation. During inference, whenever the sequence length is below 8k, RoPE is modified using the rescale factors obtained from this search procedure.

We conduct comprehensive evaluations using a range of LLMs and diverse long-context tasks to validate the efficacy of our approach. Our results indicate that LongRoPE excels at sustaining low perplexity over evaluation lengths spanning from 4k up to 2048k tokens, attains more than 90\% accuracy in passkey retrieval, and matches the accuracy of established benchmarks constructed within a 4096-token context window. Additionally, LongRoPE is compatible with all LLMs utilizing RoPE embedding. Our implementation and the LongRoPE-2048k models will be made publicly available.

\section{Non-uniformity in Positional Interpolation}
\label{sec:analysis}

\subsection{Preliminary}

Transformer models require explicit positional information, often in the form of position embedding, to represent the order of input tokens. Our work focuses on the RoPE~\cite{rope} position embedding, which is widely used in recent  LLMs. For a token at position index $n$, its corresponding RoPE encoding can be simplified as follows:

\begin{equation}
		\fontsize{7.8}{7.8} \selectfont
	\label{eq:rope}
	\left[ cos(n\theta_0), sin(n\theta_0), cos(n\theta_1), \cdot\cdot\cdot, cos(n\theta_{d/2-1}), sin(n\theta_{d/2-1}) \right]   
\end{equation}

where $d$ denotes the dimensionality of the embeddings, 
 $n\theta_i$ corresponds to the rotary angle associated with the token situated at position $n$, 
and $\theta_i=\theta^{-2i/d}$ specifies the rotational frequencies. Within RoPE, $\theta$ by default assumes the base value of 10000.

\textbf{\emph{Context window extension ratio $s$} and positional interpolation}. Here, $s$ is introduced to quantify the proportion between the extended context length $L'$ and the original context length $L$, formally: $s=\frac{L'}{L}$.

To extend the context window from $L$ to $L'$, current positional interpolation methods suggest downscaling rotation frequencies $\theta_i$ based on extension ratio $s$. Let $\beta=\theta^{2/d}$, and $\lambda$ denote the actual rescale factor related to $s$, we   unify these positional interpolation methods as follows:

\begin{equation}	
	\fontsize{7.5}{7.5} \selectfont
	\label{eq:unified_rope}
	\begin{aligned}
		\left[cos\left(\frac{n}{\lambda(\beta)^0}\right), sin \left(\frac{n}{\lambda(\beta)^0}\right), cos\left(\frac{n}{\lambda(\beta)^1}\right), \cdot\cdot\cdot,  sin \left(\frac{n}{\lambda(\beta)^{d/2-1}}\right) \right]   
	\end{aligned}
\end{equation}

\textbf{Linear positional interpolation (PI)}. The PI method~\cite{pi} advocates for using linear interpolation of positional indices within the confines of the original pre-training sequence length. When extending the sequence by a factor of $s$, it uniformly decreases the rotation angles for every position in all RoPE dimensions by $\lambda = s$. This approach compresses positional information, which reduces the model's capability to discriminate between tokens placed in close proximity. Consequently, PI's performance degrades, especially when the extension ratio becomes large.

\textbf{NTK-based interpolation and extrapolation}. The works of~\cite{ntk,dynamicntk} approach RoPE from the viewpoint of information representation, employing Neural Tangent Kernel (NTK) theory~\cite{jacot2018neural,tancik2020fourier}. To address the issue of positional crowding in PI, they propose redistributing the interpolation burden among RoPE’s dimensions. Specifically, the scaling is applied differently: dimensions with lower frequencies are scaled more aggressively, while those with higher frequencies receive less scaling. This allows for both positional interpolation and extrapolation, utilizing $\lambda=s^{i}$. The enhanced dynamic NTK~\cite{dynamicntk} adapts the extension ratio at each location depending on the current sequence length. In contrast to PI, which requires fine-tuning for extension, NTK-aware techniques permit enlarging context windows even without additional training, although these methods are generally limited to a maximum extension ratio of 4$\times$.

\textbf{YaRN}~\cite{yarn} presents a notable enhancement for positional interpolation. It separates RoPE dimensions according to their associated frequencies into three distinct categories, applying unique interpolation methods to each. Specifically, it extrapolates the high-frequency dimensions ($\lambda$=1), applies linear positional interpolation (PI) to the low-frequency group, and utilizes the NTK method for the intermediate range. A central feature of YaRN is its dimension grouping strategy; however, this relies predominantly on empirical, manual analysis. As a result, the approach may not generalize optimally to newly introduced LLM architectures.

\subsection{Study on Non-uniform Positional Interpolation}

Building on insights from NTK and YaRN, which attribute their improvements to the incorporation of non-linearity—especially through handling distinct frequency components within RoPE dimensions to enhance both interpolation and extrapolation—it becomes apparent that present methods for non-linearity are largely constrained by heuristics devised by researchers. This observation prompts two primary inquiries: (1) Does the existing approach to positional interpolation achieve the best possible results? (2) Might there exist alternative, uninvestigated forms of non-linearity?

In order to address these queries, we employ evolution search (refer to Sec.~\ref{sec:method}) to identify improved non-uniform positional interpolations specifically for LLaMA2-7B. The search process is steered by perplexity metrics, evaluated across 5 randomly selected samples drawn from the PG19~\cite{pg19} validation dataset. Our experimental investigations yield the following principal observations.

\textbf{Finding 1}: \textit{RoPE dimensions exhibit substantial non-uniformities, which are not effectively handled by current positional interpolation methods.}

We identify the optimal value of $\lambda$ for each RoPE dimension as outlined in Eq.~\ref{eq:unified_rope}. Table~\ref{tbl:exp1} presents a comparison of LLaMA2-7B’s perplexity across various methods on the PG19 and Proof-pile~\cite{proof-pile} test sets, all evaluated without fine-tuning. Our empirically optimized configuration yields notable enhancements, implying that both established linear (PI) and current non-uniform interpolation techniques (Dynamic-NTK and YaRN) are not fully optimal. Interestingly, YaRN demonstrates inferior performance relative to PI and NTK on PG19, primarily because it fails to achieve the intended context window length in the absence of fine-tuning. For instance, with an 8k context window, YaRN’s perplexity markedly increases after reaching 7k tokens.

In our investigation, we found that the rescaling factors $\lambda$ in Eq.~\ref{eq:unified_rope} are not uniform, contrasting with the constant scale $s$ used in PI, the computations employed in NTK, and YaRN's group-based rescaling. The adoption of these non-uniform factors yields marked improvements in LLaMA2’s language modeling performance (as measured by perplexity) when handling 8k and 16k context lengths, all without requiring additional fine-tuning. This improvement arises because the resultant positional embeddings more faithfully retain properties of the original RoPE, particularly in crucial dimensions, making it less challenging for the LLM to discern tokens that are near each other in the sequence.

\textbf{Finding 2}: \textit{RoPE for the initial tokens in the input sequence should be extrapolated with less interpolation.} 

We propose that for the first $\hat{n}$ tokens of input sequences, the RoPE mechanism should minimize interpolation. The rationale is that these early tokens typically receive substantial attention scores, establishing their significance within the attention layers, as previously noted in Streaming LLM~\cite{streamingllm} and LM-Infinite~\cite{han2023lminfinite}. To assess this conjecture, we expand the context window to 8k and 16k using PI and NTK methods, ensuring that the initial $\hat n$ tokens ($\hat n=0,2, ..., 256$) are not interpolated. Setting $\hat n$ to 0 reverts the system back to the standard PI and NTK configurations. Table~\ref{tbl:tokenposition} reveals two primary insights: \textbf{(1)} omitting interpolation for the beginning tokens directly results in performance enhancement; \textbf{(2)} the best value for $\hat n$ is contingent upon the target context window size.

\textbf{Finding 3}: \textit{Non-uniform positional interpolation effectively extends LLM context window in both fine-tuning and non-fine-tuning settings.} 

Although our experiments demonstrated that applying our discovered non-uniform position interpolation markedly boosts extension capability at 8k and 16k context lengths without additional training, pushing context sizes even further necessitates fine-tuning. Consequently, we perform fine-tuning of LLaMA2-7B with the optimized RoPE configuration we identified, targeting a 64k context window (details are in the Appendix). As indicated in Table~\ref{tbl:finetune}, our approach yields substantially better results than both PI and YaRN, regardless of whether LLaMA2-7B is fine-tuned. This performance gain stems from our effective employment of non-uniform positional interpolation, which reduces information degradation and delivers improved initial states for subsequent fine-tuning.

\textbf{Summary}. Our research highlights two non-uniform aspects: differential RoPE dimensions and non-uniformity in token locations. Strategically leveraging these non-uniform factors within positional interpolation leads to major gains in extending LLM context capabilities.

\section{LongRoPE}

\label{sec:method}
Building upon these observations, we propose LongRoPE. This method initially deploys an optimized search procedure that capitalizes on both forms of non-uniformity, subsequently enabling the expansion of the LLM context window to exceed 2 million tokens.

\subsection{Problem Formulation}

These two sources of non-uniformity greatly expand the possible solution space and complicate the optimization process. To manage this, we recast the optimization challenge of multidimensional non-uniform position interpolation into a search problem.

For a LLM targeting a  context window size of $L'$ and lengthy input documents $\mathbf{X}$, where each $\mathbf{x}\in\mathbf{X}$ surpasses $L'$ in token length, we denote the original rotary angle of the $i^{th}$ dimension in RoPE embedding at token position $n$ as  $\frac{n}{\beta^i}$. The optimization problem is then formulated as follows:

\begin{equation}
\begin{gathered}
		\label{eq:problem}

\underset {\mathbf{x}\in\mathbf{X}; \, |\mathbf{x}|\ge L'}{ \operatorname {arg\,min} } \,  \mathcal{L}\, \left( \text{LLM} (\text{RoPE}, \mathbf{X})\right),	\text{with}
\\
\scalebox{0.93}{
$\underset {\begin{subarray}{l} i=0,\cdot\cdot,\frac{d}{2}-1;\\ n\in[ 0,|\mathbf{x}|);\end{subarray}}{\text{RoPE}_(n)}= {\left[\cdots, \cos\left(\mathbb{I}(\hat{\lambda}_{i}, \hat{n}) \frac{n}{\beta^i}\right), \sin\left(\mathbb{I}(\hat{\lambda}_{i}, \hat{n}) \frac{n}{\beta^i}\right), \cdots\right] }$}\\
	where $\mathbb{I}(\hat{\lambda}_{i},\hat{n}) = \begin{cases}
	1 &\mbox{\;  $n<\hat{n}$}\\
	{\frac{1}{\lambda}_{i}} &\mbox{\; $n\geq\hat{n}$}
	\end{cases}$.
	\end{gathered}
\end{equation}

In this approach, we define a set of rescale factors, $\mathbb{I}(\hat{\lambda}_{i},\hat{n})$, which address the two types of non-uniformity. Here, $\hat{\lambda}_{i}$ refers to the dimensional non-uniformity of RoPE, while $\hat{n}$ indicates the positional non-uniformity for tokens. Specifically, the function $\mathbb{I}(\hat{\lambda}_{i},\hat{n})$ adjusts the rotation angle for the $i^{th}$ RoPE dimension; $\hat\lambda_{i}$ serves as the scaling coefficient and $\hat{n}$ is the threshold for token position. For token indices less than $\hat{n}$, the original rotary position embedding angle, $\frac{n}{\beta^i}$, is employed with no scaling applied. When token indices satisfy $n\ge\hat{n}$, the rescale factor $\hat\lambda_{i}$ is introduced.

Given a desired context window size of $L'$, our goal is to identify the set of optimal rescaling factors ($\mathbb{I}(\hat{\lambda}_{0},\hat{n})$, $\mathbb{I}(\hat{\lambda}_{1},\hat{n})$, ..., $\mathbb{I}(\hat{\lambda}_{i},\hat{n})$, ...) across the $1^{st}$ to $d^{th}$ dimensions of RoPE. With these rescaled RoPE parameters applied to the target $\text{LLM}$, the model can attain the lowest possible next token prediction loss, $\mathcal{L}$ (which corresponds to perplexity), on input sequences $\mathbf{X}$ consisting of $L'$ tokens.

\subsection{Searching the Non-uniform Position Interpolation}

In order to address the issue presented in Eq.~\ref{eq:problem}, we propose a straightforward and powerful approach that systematically explores optimal RoPE rescale factors, enabling comprehensive utilization of the multidimensional variations present in position embeddings.

\textbf{Search space}. We construct an expansive search space that captures both types of non-uniformity. Table~\ref{tbl:searchspace} provides an overview of this search space. In detail, we enable searching for a unique rescale factor along each dimension of the RoPE embedding. To streamline the search space, we parameterize it using $\lambda_i$ and $\hat{n}$ in lieu of directly searching over $\mathbb{I}(\hat{\lambda}_{i},\hat{n})$, with the substitution $\hat{\lambda}_{i}=1/\lambda_i$. As depicted in Table~\ref{tbl:searchspace}, $\lambda_i$ may be chosen from a lower bound of 1.0 (corresponding to pure extrapolation) up to a maximum of $s\times1.25$ (enabling broader interpolation than permitted by PI), incremented in steps of 0.01, where $s$ refers to the target expansion ratio for the context window.

The parameter $\hat{n}$ specifies how many of the first token positions keep their original RoPE embeddings, foregoing position interpolation. In practice, we permit $\hat{n}$ to vary over the set \{0, 1, 2, 4, 8, 12, 16, 20, 24, 28, 32, 64, 128, 256\}. If $\hat{n}=0$, every token position applies the identified rescale factors.

\textbf{Evolution-based search}. The search space defined in Table~\ref{tbl:searchspace} encompasses a vast array of possible positional interpolation configurations, making comprehensive exploration highly challenging. As an example, setting $s=4\times$ results in $400^{128/2}\times14=4\times10^{167}$ alternatives. This number grows exponentially with increasing extension ratios. To mitigate these difficulties, we employ an evolution search strategy~\cite{guo2020single} complemented by two specialized optimization methods that substantially enhance search efficiency. The complete search workflow is detailed in Algorithm~\ref{alg:search}.

\textit{Optimized initial population generation}. Rather than solely relying on random initialization for the population of $P$ rescale factors, we incorporate the three RoPE rescale factors associated with PI, NTK, and YaRN directly as members of the initial population. The remaining $P$-3 individuals are then created by applying random mutations, each occurring with probability $p$, to the values of these three rescale factors.

\textit{Monotonically non-decreasing constraint}. Following the creation of the initial population, we assess the LLM perplexity of every candidate. This is accomplished by applying the relevant RoPE rescale factors to the target LLM and evaluating the perplexity on input $\mathbf{X}$. The highest-performing $k$ individuals are chosen as parents for subsequent evolutionary steps. Nonetheless, the sheer size of the search space means that indiscriminate mutation and crossover can lead to suboptimal candidates, resulting in excessive and redundant perplexity evaluations. This inefficiency is exacerbated when $L'$ is large, since each perplexity inference demands significant computational resources.

To mitigate this issue, we enforce a constraint that requires the sampled RoPE rescaled factors to be monotonically non-decreasing: $\lambda_i \le \lambda_{i+1}$. Only those RoPE parameters conforming to this monotonicity are considered for perplexity evaluation in the LLM, thereby substantially streamlining the search process. In detail, we stipulate that $\lambda_i$ must increase with respect to the RoPE dimension ($i=0,\ldots,63$). This monotonicity across dimensions is motivated by NTK theory~\cite{jacot2018neural,tancik2020fourier,ntk}, which posits that dimensions associated with higher frequencies—i.e., lower indices—should employ less interpolation (smaller $\lambda_i$), while those corresponding to lower frequencies—higher indices—are capable of greater interpolation (larger $\lambda_i$).

\begin{algorithm}[t]
	\small
	\caption{The search algorithm}
	\textbf{Input:} target LLM, input samples $\mathbf{X}$, population size $P$, mutation size $N_1$, crossover size $N_2$, maximum iterations $\mathcal{T}$, mutation probability $p$\\

	\begin{algorithmic}[1]
		\label{alg:search}
		\STATE Top-k $\gets \phi$;	
		\STATE $\text{P}_0 \gets$ \textit{Initialize\_population\_with\_optimization}($P$, $\mathbf{X}$, $p$);
		\FOR{$i=1$ to $\mathcal{T}$}
		\STATE \textit{Compute\_perplexity}(LLM, $\text{P}_{i-1}$, $\mathbf{X}$);
		\STATE Top-k $\gets$ \textit{Update\_Topk}(Top-k, $\text{P}_{i-1}$);
		\STATE $\text{P}_{mutation} \gets$ \textit{Mutation\_with\_mono\_constraint}(Top-k, $N_1$, $p$);
		\STATE $\text{P}_{crossover} \gets$ \textit{Crossover\_with\_mono\_constraint}(Top-k, $N_2$);
		\STATE $\text{P}_i \gets \text{P}_{mutation} \cup \text{P}_{crossover} \cup$ Top-k;
		\ENDFOR
		\STATE Output the individual from Top-k that achieves the minimum perplexity;
	\end{algorithmic}
\end{algorithm}

\textbf{8$\times$ extension without fine-tuning}. Leveraging evolutionary search, we systematically determine optimal, non-uniform RoPE rescaling factors that protect critical dimensions and positional information, thereby mitigating information degradation due to interpolation. As illustrated in Fig.\ref{fig:non-ft}, this approach enables the context window of LLaMA2 to be expanded from 4k to 32k tokens without the need for fine-tuning. In comparison, prior techniques like PI, as well as non-uniform NTK and YaRN, exhibit significant increases in perplexity once the context window is doubled.

\subsection{Extending LLM Context Window to 2048K}
\label{sec:extendto2048k}

\textbf{Progressive extension to 2048k}. In this section, we detail our approach for increasing the context window size of existing pre-trained LLMs from the standard 4k tokens up to beyond 2048k tokens. Our approach with non-uniform positional interpolation enables up to 8$\times$ extension without requiring model fine-tuning. However, in cases where much larger expansions (such as 512$\times$) are needed, fine-tuning becomes indispensable. A common strategy entails identifying suitable RoPE rescaling coefficients for the 2048k target window before initiating fine-tuning. Despite this, several obstacles remain, chiefly the immense computational cost involved in such large-scale training. Furthermore, we have observed in practice that fine-tuning LLMs with extremely high extension ratios is particularly difficult to accomplish effectively (see Appendix).

Luckily, LongRoPE demonstrates robust performance on both base and fine-tuned extended LLMs. As a result, we propose a streamlined, incremental approach that reaches the desired 2048k sequence length using only 1k fine-tuning iterations, limited to a maximum training length of 256k.

$\Diamond$ \textit{LongRoPE search for pre-trained LLM expansion to 256k}. Using LLaMA2 as a case study, we perform searches aimed at increasing the context window to 128k and 256k tokens. Here, the context length is extended by factors of 32$\times$ and 64$\times$ relative to the original, respectively.

$\Diamond$ \textit{Fine-tuning to 256k}. After pre-training, we adapt the LLM to handle a context window of 256k. In this approach, we begin by fine-tuning LLaMA2 over 400 steps with RoPE rescaled factors set for 128k. Subsequently, we update the RoPE rescaled factors to 256k on the resulting checkpoint and carry out a further 600 fine-tuning steps. Compared to directly fine-tuning for the full 256k context window, this strategy demonstrates improved efficiency.

$\Diamond$ \textit{Scaling fine-tuned extended LLMs to 2048k using LongRoPE search}. In the final step, we apply a subsequent search procedure to the previously fine-tuned LLM at the 256k length. This approach allows us to achieve a context window spanning 2048k tokens, all without requiring additional fine-tuning. This yields a maximum extension ratio of $512\times$.

\textbf{Performance restoration within limited context windows}. Upon increasing the context window to a vast 2048k length, a decline in performance is observed for sequences contained within the initial window size. This phenomenon has been documented in positional interpolation research~\cite{pi}, which demonstrates that when higher-dimensional position embeddings are compressed into a restricted region within the standard context window, language model effectiveness is adversely impacted. Using an extension ratio of 512$\times$ causes the original 4k context window positions to be densely packed, exacerbating this issue.

To address this issue, an additional evolution search is conducted on the extended LLM to fine-tune the RoPE rescale factors specifically for shorter context lengths such as 4k and 8k. Since these shorter lengths require minimal positional interpolation, the upper bound for the searched $\lambda$ is lowered. In inference, the LLM adaptively modifies the relevant RoPE rescale factors.

\section{LongRoPE}

\label{sec:method}
Inspired by these observations, we propose LongRoPE, which initially implements a streamlined search method to harness both types of non-uniformity. This technique is then leveraged to push the LLM context window to lengths exceeding 2 million tokens.

\subsection{Problem Formulation}

These two irregularities contribute to an expansive range of possible solutions, thereby increasing the intricacy of the optimization process. To mitigate this challenge, we reformulate the multidimensional optimization of non-uniform position interpolation into a search-based problem.

For a LLM targeting a  context window size of $L'$ and lengthy input documents $\mathbf{X}$, where each $\mathbf{x}\in\mathbf{X}$ surpasses $L'$ in token length, we denote the original rotary angle of the $i^{th}$ dimension in RoPE embedding at token position $n$ as  $\frac{n}{\beta^i}$. The optimization problem is then formulated as follows:

\begin{equation}
\begin{gathered}
		\label{eq:problem}

\underset {\mathbf{x}\in\mathbf{X}; \, |\mathbf{x}|\ge L'}{ \operatorname {arg\,min} } \,  \mathcal{L}\, \left( \text{LLM} (\text{RoPE}, \mathbf{X})\right),\quad \text{where} 
\\
\scalebox{0.93}{
$\underset {\begin{subarray}{l} i=0,\ldots,\frac{d}{2}-1;\\ n\in[ 0,|\mathbf{x}|);\end{subarray}}{\text{RoPE}_(n)}= {\left[\ldots, \cos\left(\mathbb{I}(\hat{\lambda}_{i}, \hat{n})\cdot\frac{n}{\beta^i}\right), \sin\left(\mathbb{I}(\hat{\lambda}_{i}, \hat{n})\cdot\frac{n}{\beta^i}\right), \ldots \right] }$}\\
	\text{with} \quad \mathbb{I}(\hat{\lambda}_{i},\hat{n})= \begin{cases}
	1 &\text{if $n<\hat{n}$}\\
{\frac{1}{\lambda}_{i}} &\text{if $n\geq\hat{n}$}
\end{cases}
	\end{gathered}
\end{equation}

In this context, a set of scaling parameters $\mathbb{I}(\hat{\lambda}_{i},\hat{n})$ is incorporated to address two types of non-uniformity. Here, $\hat{\lambda_i}$ represents the irregularity within RoPE dimensions, while $\hat{n}$ reflects token positional non-uniformity. More concretely, $\mathbb{I}(\hat{\lambda}_{i},\hat{n})$ modifies the rotational angle associated with the $i^{th}$ dimension of RoPE by utilizing $\hat\lambda_{i}$ as the scaling factor and $\hat{n}$ as the threshold for token position. When considering the first $\hat{n}$-1 token indices, the scaling effect of $\hat\lambda_{i}$ is disregarded and the unadjusted RoPE angle $\frac{n}{\beta^i}$ is applied. For token positions satisfying $n\ge\hat{n}$, the scaling factor then becomes operative.

With a desired context window length of $L'$, our task is to identify the most suitable rescale factors ($\mathbb{I}(\hat{\lambda}_{0},\hat{n})$, $\mathbb{I}(\hat{\lambda}_{1},\hat{n})$, ...$\mathbb{I}(\hat{\lambda}_{i},\hat{n})$...) corresponding to each RoPE dimension, from the $1^{st}$ up to the $d^{th}$. The goal is to modify the target $\text{LLM}$ with these adapted $\text{RoPE}$ scalings in such a way that it minimizes next token prediction loss, $\mathcal{L}$ (which equates to perplexity), over input sequences $\mathbf{X}$ of token length $L'$.

\subsection{Searching the Non-uniform Position Interpolation}

In addressing the challenge posed by Eq.~\ref{eq:problem}, we present a straightforward and potent approach that identifies the best RoPE rescale factors, enabling comprehensive utilization of the multidimensional asymmetries inherent in position embedding.

\textbf{Search space}. We establish an extensive search space to encompass both sources of non-uniformity. Table~\ref{tbl:searchspace} provides details about this configuration. In particular, each dimension within the RoPE embedding is assigned a customizable rescale factor, enabling more granular optimization. For simplicity in constructing the search space, the optimization is performed over $\lambda_i$ and $\hat{n}$ rather than directly tuning $\mathbb{I}(\hat{\lambda}_{i},\hat{n})$, given that $\hat{\lambda}_{i}=1/\lambda_i$. As depicted in Table~\ref{tbl:searchspace}, the range of possible values for $\lambda_i$ extends from a lower bound of 1.0 (corresponding to pure extrapolation) up to $s\times1.25$ (representing interpolation more aggressive than PI), incremented in steps of 0.01, where $s$ denotes the factor by which the target context window is to be extended.

$\hat{n}$ determines how many starting token positions are preserved with their original positional encoding, meaning they do not undergo position interpolation and instead retain the initial RoPE embedding. In practice, we set $\hat{n}$ to explore values within \{0, 1, 2, 4, 8, 12, 16, 20, 24, 28, 32, 64, 128, 256\}. When $\hat{n}=0$, rescale factors found through search are applied uniformly across all token positions.

\textbf{Evolution-based search}. As detailed in Table~\ref{tbl:searchspace}, our search space comprises a vast array of positional interpolation configurations, making exhaustive exploration computationally infeasible. For instance, considering an extension ratio of $s=4\times$, the possible combinations amount to $400^{128/2}\times14$=4$\times10^{167}$, demonstrating a rapid exponential growth in the search domain as the extension ratio increases. To efficiently navigate this immense landscape, we employ evolution search~\cite{guo2020single} and implement two key optimization strategies that significantly enhance the search efficiency. The entire search workflow is depicted in Algorithm~\ref{alg:search}.

\textit{Optimized initial population generation}. Rather than relying solely on random initialization for the population of $P$ rescale factors, our approach incorporates the three RoPE rescale factors associated with PI, NTK, and YaRN directly as members of the initial population. For the other $P$-3 individuals, we introduce random mutations to these three rescale factors with a likelihood of $p$.

\textit{Monotonically non-decreasing constraint}. Following the creation of the initial population, each individual's LLM perplexity is evaluated. In this step, we adjust the target LLM using the associated RoPE rescale factors, then determine the perplexity score for the input $\mathbf{X}$. The highest-scoring $k$ individuals are selected as parents for subsequent evolutionary steps. The enormous search space, however, often results in random mutation and crossover operations encountering low-quality candidates, which translates to superfluous perplexity evaluations. This inefficiency is further exacerbated for large values of $L'$, as each perplexity inference entails considerable computational cost.

To mitigate this issue, we enforce a non-decreasing monotonicity restriction on the sampled RoPE rescale coefficients: $\lambda_i\le \lambda_{i+1}$.
Only those RoPE configurations adhering to this monotonic criterion are implemented within the LLM during perplexity assessment, thereby markedly curtailing search overhead. More concretely, we stipulate that $\lambda_i$ should grow monotonically with respect to RoPE dimension (that is, $i=0,\ldots,63$). This requirement of dimensional monotonicity is informed by NTK theory~\cite{jacot2018neural,tancik2020fourier,ntk}, which posits that low-indexed dimensions—associated with higher-frequency components—necessitate less interpolation (hence, smaller $\lambda_i$ values), whereas those dimensions at higher indices—corresponding to lower-frequency elements—can accommodate greater interpolation (thus, larger $\lambda_i$ values).

\begin{algorithm}[t]
	\small
	\caption{The search algorithm}
	\textbf{Input:} target LLM, input samples $\mathbf{X}$,   population size $P$, mutation size $N_1$, crossover size $N_2$, max iterations $\mathcal{T}$, mutate probability $p$\\

	\begin{algorithmic}[1]
		\label{alg:search}
		\STATE Top-k $\gets$ $\phi$;	
		\STATE $\text{P}_0 \gets$ \textit{Initialize\_population\_with\_optimization}($P$, $\mathbf{X}$, $p$); 
		\FOR{$i = 1$ \TO $\mathcal{T}$}
		\STATE \textit{Compute\_perplexity}(LLM, $\text{P}_{i-1}$, $\mathbf{X}$);
		\STATE Top-k $\gets$ \textit{Update\_Topk}(Top-k, $\text{P}_{i-1}$);
		\STATE $\text{P}_{mutation} \gets$ \textit{Mutation\_with\_mono\_constraint}(Top-k, $N_1$, $p$); 
		\STATE $\text{P}_{crossover} \gets$ \textit{Crossover\_with\_mono\_constraint}(Top-k, $N_2$);
		\STATE $\text{P}_i \gets$ $\text{P}_{mutation} \cup \text{P}_{crossover} \cup$ Top-k;
		\ENDFOR
		\STATE Return the individual with minimum perplexity in Top-k;
	\end{algorithmic}
\end{algorithm}

\textbf{8$\times$ extension without fine-tuning}. Utilizing evolutionary search, we efficiently determine optimal non-uniform RoPE rescale factors that retain crucial dimensions and positional information, thereby reducing information loss associated with interpolation. As illustrated in Fig.\ref{fig:non-ft}, our approach successfully expands LLaMA2’s context window from 4k to 32k in the absence of fine-tuning. By comparison, current techniques like PI as well as non-uniform NTK and YaRN result in a sharp increase in perplexity following a 2$\times$ context window extension.

\subsection{Extending LLM Context Window to 2048K}
\label{sec:extendto2048k}

\textbf{Progressive extension to 2048k}. Here, we present our approach for scaling the context window of pre-trained LLMs beyond the conventional 4k limit, achieving capacities upwards of 2048k tokens. Our findings indicate that applying non-uniform positional interpolation enables an 8$\times$ increase without necessitating model fine-tuning. However, when significantly larger expansions are desired (such as up to 512$\times$), fine-tuning becomes indispensable. A potential strategy involves identifying suitable RoPE rescaling factors for the intended 2048k context window, followed by fine-tuning. Nevertheless, this process is impeded by extremely high computational demands. Furthermore, our observations suggest that achieving effective fine-tuning becomes increasingly difficult when attempting to support such large context expansions (refer to Appendix).

Fortunately, LongRoPE demonstrates strong performance with both the base model and the extended LLM post fine-tuning. As a result, we propose a resource-efficient, incremental approach that reaches the desired 2048k context window using only 1k fine-tuning iterations, and confines training to sequences of up to 256k length.

$\Diamond$ \textit{LongRoPE search enables 256k extension of pre-trained LLMs}. Using LLaMA2 as a case study, we perform a search aimed at increasing the context window to 128k and 256k tokens. At this phase, the corresponding extension factors are 32$\times$ and 64$\times$, accordingly.

$\Diamond$ \textit{Fine-tuning to 256k}. Following pre-training, we proceed to adapt the LLM for a 256k context window via fine-tuning. The procedure begins with 400 fine-tuning steps of LLaMA2 utilizing RoPE rescaled factors corresponding to 128k. After this initial phase, the RoPE rescaled factors are updated to 256k on the resulting checkpoint, and an extra 600 fine-tuning steps are performed. This two-stage approach has demonstrated greater efficiency compared to fine-tuning directly to 256k.

$\Diamond$ \textit{Scaling fine-tuned extended LLM to 2048k using LongRoPE search}. In the concluding step, an additional search is conducted on the LLM that has already been fine-tuned for a 256k context length. This process enables expansion to a substantially larger window of 2048k, achieved without additional fine-tuning. The overall scale-up factor reaches $512\times$.

\textbf{Restoration of performance in shorter context windows}. Upon expanding the context window to a sizable 2048k tokens, we observe a degradation in model performance within the initial context window. This phenomenon has been previously identified with positional interpolation~\cite{pi}, as the method compresses positional embeddings within the original window into a limited subspace of higher-dimensional representation. Such compression adversely impacts the effectiveness of the language model. Notably, with an extension factor of 512$\times$, token positions that previously spanned the initial 4k context window are now densely packed, exacerbating this issue.

To address this issue, we conduct an additional evolutionary search on the extended LLM to fine-tune the RoPE rescale factors specifically for shorter context lengths, such as 4k and 8k. For these shorter ranges, we decrease the upper bound of the searched $\lambda$ values, reflecting the diminished need for positional interpolation. At inference time, the LLM adapts the relevant RoPE rescale factors dynamically.

 \section{Experiments}

\label{sec:exp}
\subsection{Setup}
\textbf{Evaluation Tasks and models}. Our experiments utilize LongRoPE integrated with both LLaMA2-7B and Mistral-7B. Model capabilities are assessed in three main areas: (1) examining perplexity on lengthy documents using extended-context LLMs; (2) the Passkey retrieval task, designed to test the capacity to extract a target passkey embedded within a large amount of distracting content; and (3) performance on established LLM benchmarks, restricted to a 4096-token context window.

\textbf{Fine-tuning}. For LLaMA2, we adopt a linear learning rate schedule starting at 2e-5 and maintain a global batch size of 32 throughout the training process. The model undergoes 400 fine-tuning steps using the Redpajama~\cite{redpajama} dataset, which is divided into 128k-length segments with BOS and EOS tokens marking segment boundaries. After this initial phase, an additional 600 training steps are executed, building on the resulting checkpoint, to extend the context window to 256k tokens. Fine-tuning at the 128k context length leverages 8 A100 GPUs via the distributed system~\cite{cube}; scaling to 256k necessitates 16 A100 GPUs. For Mistral, training employs a fixed learning rate of 1e-6 and a global batch size set to 64. The fine-tuning procedure for both 128k and 256k context windows adheres to the YaRN protocol~\cite{yarn}, consisting of 400 steps on Together Computer's Long-Data Collections~\cite{mistraldata} with sequence lengths of 16k. All Mistral models are trained using 4 A100 GPUs.

\textbf{Search}. For experiments with a target window size up to 256k, the following parameters are applied: $P$=64, $N_1$=$N_2$=16, $p$=0.3, and $\mathcal{T}$=40. In each iteration, the top-32 candidates are chosen for subsequent mutation or crossover steps. Perplexity is assessed using 5 randomly selected samples from the PG19 validation set, each sample meeting at least the specified context length. When operating with windows larger than 512k, both the population size and the numbers for mutation and crossover are reduced by half. In this scenario, perplexity is evaluated on 3 randomly chosen Pile-Books3~\cite{pile} validation set samples.

\textbf{Baselines}. In order to achieve a context window of 2048k, we fine-tuned models that initially supported 128k and 256k context lengths. This results in LongRoPE-2048k (ft=128k) and LongRoPE-2048k (ft=256k) for the LLaMA2 and Mistral architectures, respectively. We evaluate these four models against leading context window extension baselines, focusing on open-source LLMs that were fine-tuned following positional interpolation methods such as PI, NTK, and YaRN. Among these are Together-32k~\cite{together}, Code LLaMA~\cite{codellama}, LongLoRA-full-FT-100k~\cite{longlora}, YaRN-LLaMA, and YaRN-Mistral~\cite{yarn}.

\subsection{Main Results}

\label{sec:mainresult}
\textbf{Language modeling for lengthy sequences up to 256k}. Our analysis starts with a comparison against leading extended LLM models when evaluated on sequences with a maximum length of 256k tokens. To illustrate the robustness of our approach, we utilize two distinct datasets: the Proof-pile~\cite{pg19} and the PG19~\cite{pile} test splits. Perplexity is measured at multiple context sizes by applying a sliding window with a width of 256 tokens. For the PG19 dataset, the evaluation leverages the complete test set, comprising 100 documents. For Proof-pile, following the protocol outlined in YaRN~\cite{yarn}, we randomly select 10 samples, ensuring that each sample contains no fewer than 128k tokens.

Table~\ref{tbl:proof-pile} and Table~\ref{tbl:pg19} present perplexity measurements for LLaMA2 and Mistral augmented with various interpolation strategies, evaluated on Proof-pile and PG19 datasets, respectively. Two main insights emerge from the results: \textbf{(1)} The interpolated models consistently demonstrate lower perplexity as the evaluation length increases from 4k up to 256k, indicating enhanced capacity to utilize extended contextual information. \textbf{(2)} Remarkably, despite operating with context windows that are up to 16$\times$ larger—an especially demanding scenario for sustaining short-context performance—the LongRoPE-2048k configurations surpass leading baseline models in the 256k context regime.

\textbf{Language modeling for sequences exceeding 2000k tokens}. To assess performance on ultra-long textual data, we utilize the Books3~\cite{pile} dataset. For efficient evaluation, we randomly choose 20 books, all surpassing 2048k tokens, and employ a sliding window approach with a span of 256k tokens.

Table~\ref{tbl:books3} demonstrates that LongRoPE effectively increases the context window of both LLaMA2-7B and Mistral-7B to 2048k tokens, while maintaining perplexity levels on short ranges (8k–128k) that are either similar to or better than the baseline models. Distinct differences emerge between LLaMA2 and Mistral at the 2048k length: Mistral shows superior performance at shorter context sizes, but its perplexity rises above 7 when the window goes beyond 256k. LLaMA2, in contrast, exhibits expected behavior with a consistent decrease in perplexity as the context length expands, though slight increases are noticed at 1024k and 2048k tokens. Notably, for LLaMA2, fine-tuning LongRoPE-2048k at a 256k window yields better results than at 128k, attributed to the lower extension ratio (8$\times$ as opposed to 16$\times$). Conversely, Mistral achieves better results when fine-tuned at a window size of 128k. This pattern arises from the fact that both 128k and 256k fine-tuning in Mistral leverage a training length of 16k, as prescribed by YaRN's protocol, which limits Mistral's capacity to scale its context window during subsequent fine-tuning.

\textbf{Passkey retrieval}. Next, we examine how the effective context window influences performance in generation-based scenarios. For this purpose, we utilize a synthetic passkey retrieval benchmark as outlined by ~\cite{passkey}. This benchmark requires the model to locate a randomly chosen five-digit passkey embedded within an extended text. Additional information about the prompt format can be found in the appendix. We conduct ten repetitions of the passkey retrieval experiment, randomly inserting the passkey at locations drawn from a uniform distribution over the full evaluation context window.

As illustrated in Fig.~\ref{fig:passkey}, the retrieval accuracy of baseline models declines sharply, reaching zero for sequence lengths beyond 128k tokens. Conversely, LongRoPE-LLaMA2-2048k (ft=256k) demonstrates impressive performance in the demanding scenario of extracting a passkey from sequences containing up to millions of tokens, sustaining retrieval accuracy rates of at least 90\% between 4k and 2048k tokens. Similarly, LongRoPE-Mistral-2048k (ft=128k) preserves perfect accuracy up to 1800k tokens before decreasing to 60\% at 2048k, a result consistent with observations from Table~\ref{tbl:books3}, where a modest perplexity increase is reported at the maximum length.

\textbf{Baseline evaluations in the native context length}. We assess LongRoPE-2048k models within their original context window employing the Hugging Face Open LLM Leaderboard~\cite{huggingfaceboard} across both zero-shot and few-shot conditions. Specifically, we utilize the ARC-Challenge~\cite{allenai:arc} in a 25-shot format, HellaSwag~\cite{zellers2019hellaswag} with 10 shots, MMLU~\cite{mmlu} using 5 shots, and TruthfulQA~\cite{lin2021truthfulqa} in the zero-shot setting.

As demonstrated in Table~\ref{tbl:commonsense}, our models deliver results similar to those on the initial benchmark, which was tailored for shorter context lengths, and surpass the original Mistral on TruthfulQA by a margin of +0.5\%. While LongRoPE-LLaMA2-2048k, trained with a 256k context length, experiences somewhat greater performance decline, its scores are still largely consistent with expectations across the majority of tasks.

\subsection{Ablation Results}

\textbf{Evaluating the second stage of positional interpolation}. Within the framework of our progressive extension approach, we implement a secondary non-uniform positional interpolation through our search algorithm on extended LLMs that have undergone fine-tuning. To assess its impact, we perform a series of experiments using the fine-tuned LLaMA2-256k model as a baseline. We enlarge the model to 512k, 1024k, and 2048k sequences, applying PI and YaRN for comparison. The results, summarized in Table~\ref{tbl:secondsearch}, indicate that our method of non-uniform positional interpolation maintains a stable perplexity value across these extensions. By contrast, both PI and YaRN exhibit a marked rise in perplexity as the extension ratio increases.

\textbf{Recovery performance with reduced context lengths.} To address the degradation in performance at smaller context windows, we recalibrate the RoPE parameters for LongRoPE-2048k using our optimization strategy. In particular, by lowering the upper bound of scale factors in the search process, we intentionally limit interpolation effects for 4k and 8k contexts. Table~\ref{tbl:shortperformance} presents a perplexity evaluation for LongRoPE-LLaMA2-2048k on Proof-pile with 4k and 8k window sizes, accompanied by the mean benchmark accuracy for LLMs. The findings indicate a marked enhancement in performance for shorter contexts.

\textbf{Analysis on the two forms of non-uniformities}. To disentangle the influence of each non-uniformity, we conducted ablation studies to assess their contributions to model performance. Our experimental setup included two scenarios: (i) scaling LLaMA2-7B to shorter sequence lengths of 16k and 32k via several approaches—PI, optimizing RoPE dimensions alone, and optimizing both types of non-uniformities; (ii) extending a previously fine-tuned LLaMA2 at 256k tokens up to 2048k using the analogous procedures. Perplexity measurements were performed in the absence of additional fine-tuning. Table~\ref{tbl:searchdimension} demonstrates that tuning non-uniformity in RoPE dimensions yields a notable reduction in perplexity over PI’s linear interpolation. Adjusting token position non-uniformity produces measurable gains for shorter contexts (16k and 32k) but fails to offer similar benefits at 2048k, likely due to the extraordinarily long sequence length. Retaining only initial tokens without interpolation does not provide further advantage in this setting; we recognize this limitation and propose further investigation.

\section{Related Works}

Beyond position interpolation techniques, this section explores related studies employing alternative strategies.

\textbf{Retrieval-based} approaches utilize an external memory system to store extensive contextual information and implement retrieval components to access pertinent documents during inference~\cite{longllama,wang2023augmenting,borgeaud2022improving}. Such approaches often require significant architectural changes to LLM models. In contrast, our methodology adopts a more streamlined approach, involving only minimal adjustments to positional embeddings. Moreover, our method extends its applicability to a broader spectrum of long context tasks, such as long-form summarization and few-shot learning, not limited to retrieval settings.

\textbf{Attention-based context window extensions}. In addition to positional embedding interpolation, certain studies extend input context within the limits of the original LLM window by altering attention mechanisms~\cite{han2023lminfinite, streamingllm,parallelwindow}. Their primary approach is to address the problem of excessive attention computation triggered by the introduction of new positions, employing innovative attention masking strategies. These approaches serve as complements to methods based on positional interpolation.

\textbf{Fine-tuning based approaches} are dedicated to optimizing the process of adapting pre-trained LLMs to accommodate extended sequence lengths by adjusting their position embeddings. Approaches such as Code LLaMA~\cite{codellama}, LLaMA2 Long~\cite{xiong2023effective}, and ScaledRoPE~\cite{liu2023scaling} employ fine-tuning techniques with an enlarged base parameter in RoPE to target specific maximum context lengths. In contrast, our proposed method demonstrates adaptability across different desired sequence lengths, scaling up to over 2M tokens. Recently, as scaling up fine-tuning to lengths greater than 128k tokens imposes considerable demands on GPU capabilities, solutions like LongLoRA~\cite{longlora} and PoSE~\cite{zhu2023pose} have been developed to address resource efficiency. Notably, our approach operates independently of such efficient fine-tuning strategies.

\section{Conclusion}

This paper introduces LongRoPE, a technique designed to dramatically increase the context capacity of LLMs to 2048k tokens, while preserving their performance in the original, shorter context ranges. Our approach leverages two distinct forms of non-uniformity present in RoPE positional embeddings, utilizing a highly efficient evolutionary search procedure. This results in two main advantages: it yields advantageous initialization for fine-tuning processes and allows an extension of the context window by 8$\times$ without the need for fine-tuning. Furthermore, we describe a stepwise expansion methodology that employs LLMs fine-tuned on 256k-length sequences, progressively scaling the context window to 2048k without incurring additional fine-tuning costs. Comprehensive experimental studies confirm the effectiveness of LongRoPE. We anticipate that LongRoPE-2048k will open up numerous possibilities for applications requiring long context windows and foster new lines of investigation in the field.

\section*{Broader Impacts} The research described in this paper aims to further the progress of Machine Learning as a field. While our findings may have a variety of societal implications, we do not identify any that require particular emphasis at this time.

	\balance

\end{document}